% GENERATED FILE. Edit canonical lessons, then run npm run generate:books.
% canonical-source-hash: fnv1a64:11974e17f4d2145f
% canonical-lessons: HI-C367-sarahna, HI-C367-tarif, HI-C367-qasur, HI-C367-alocna, HI-C367-jhelna

\chapter{Praise and Blame}
\label{ch:hi-praise-and-blame}
\hlchaptermodality{367}

\begin{chapteropening}
\textbf{By the end of this chapter:} \emph{I can say who was praised and who was blamed.}

Complete the last lesson of chapter 367: I can say who was praised and who was blamed.
\end{chapteropening}

\section[\texorpdfstring{\hi{सराहना} (sarāhnā)}{सराहना (sarāhnā)}]{\hi{सराहना} (sarāhnā) — to commend, to praise}

\label{lesson:HI-C367-sarahna}



\begin{quote}
Before the new one: say the Hindi for to mislead, then the Hindi for compelled.
\end{quote}

\subsection*{You'll want to know: \hi{सराहना}}
\textbf{\hi{सराहना}} — \emph{sarāhnā} — \textquotedblleft{}to commend, to praise\textquotedblright{}.

\begin{quote}
\hi{सबने} \hi{उसके} \hi{काम} \hi{को} \hi{सराहा।}
\end{quote}

\begin{quote}
\emph{sabne uske kām ko sarāhā.}
\end{quote}

\begin{quote}
Everyone praised his work.
\end{quote}

Past, with \hi{ने} (he or she …): \hi{उसने} \hi{मेरे} \hi{काम} \hi{को} \hi{सराहा}, \emph{usne mere kām ko sarāhā}.

\begin{cousinweb}
\textbf{\hi{सराहना}} (\emph{sarāhnā}) is a tadbhava, inherited from Sanskrit \emph{ślāghate}, \textquotedblleft{}praises\textquotedblright{} (CDIAL 12734).
\end{cousinweb}

\begin{grammarlens}[title={What you've built}]
One more word for giving reasons and opinions.
\end{grammarlens}

\subsection*{Your turn}
\begin{itemize}
\raggedright
  \item \emph{Say it:} \emph{sarāhnā}
  \item \emph{Say it:} \emph{sarāhnā}, once more
  \item \emph{Say it:} \emph{sabne uske kām ko sarāhā.} — Everyone praised his work.
  \item \emph{Recall:} say the Hindi for to mislead, then the Hindi for compelled, then say \emph{sarāhnā} again
\end{itemize}

\begin{tcolorbox}[breakable,colback=teal!4,colframe=teal!35!black,title={Before you move on}]
What does \hi{सराहना} mean? (\textquotedblleft{}To commend, to praise\textquotedblright{}.) Say it once more.
\end{tcolorbox}
\section[\texorpdfstring{\hi{तारीफ़} (tārīf)}{तारीफ़ (tārīf)}]{\hi{तारीफ़} (tārīf) — praise}

\label{lesson:HI-C367-tarif}



\begin{quote}
Before the new one: say the Hindi for compelled, then the Hindi for to commend.
\end{quote}

\subsection*{You'll want to know: \hi{तारीफ़}}
\textbf{\hi{तारीफ़}} — \emph{tārīf} — \textquotedblleft{}praise\textquotedblright{}.

\begin{quote}
\hi{शिक्षक} \hi{ने} \hi{मीरा} \hi{की} \hi{तारीफ़} \hi{की।}
\end{quote}

\begin{quote}
\emph{śikṣak ne Mīrā kī tārīf kī.}
\end{quote}

\begin{quote}
The teacher praised Mira.
\end{quote}

Feminine.

\begin{cousinweb}
\textbf{\hi{तारीफ़}} (\emph{tārīf}) is usually traced to Arabic \emph{taʿrīf}, \textquotedblleft{}making known, description\textquotedblright{}, which reached Hindi through Persian. The sense of praise is a Persian and Hindi one.
\end{cousinweb}

\begin{grammarlens}[title={What you've built}]
One more word for giving reasons and opinions.
\end{grammarlens}

\subsection*{Your turn}
\begin{itemize}
\raggedright
  \item \emph{Say it:} \emph{tārīf}
  \item \emph{Say it:} \emph{tārīf}, once more
  \item \emph{Say it:} \emph{śikṣak ne Mīrā kī tārīf kī.} — The teacher praised Mira.
  \item \emph{Recall:} say the Hindi for compelled, then the Hindi for to commend, then say \emph{tārīf} again
\end{itemize}

\begin{tcolorbox}[breakable,colback=teal!4,colframe=teal!35!black,title={Before you move on}]
What does \hi{तारीफ़} mean? (\textquotedblleft{}Praise\textquotedblright{}.) Say it once more.
\end{tcolorbox}
\section[\texorpdfstring{\hi{क़सूर} (qasūr)}{क़सूर (qasūr)}]{\hi{क़सूर} (qasūr) — fault, blame}

\label{lesson:HI-C367-qasur}



\begin{quote}
Before the new one: say the Hindi for to commend, then the Hindi for praise.
\end{quote}

\subsection*{You'll want to know: \hi{क़सूर}}
\textbf{\hi{क़सूर}} — \emph{qasūr} — \textquotedblleft{}fault, blame\textquotedblright{}.

\begin{quote}
\hi{यह} \hi{तुम्हारा} \hi{क़सूर} \hi{नहीं} \hi{है।}
\end{quote}

\begin{quote}
\emph{yah tumhārā qasūr nahīṃ hai.}
\end{quote}

\begin{quote}
It isn't your fault.
\end{quote}

Masculine.

\begin{cousinweb}
\textbf{\hi{क़सूर}} (\emph{qasūr}) is usually traced to Arabic \emph{quṣūr}, \textquotedblleft{}a shortcoming\textquotedblright{}, which reached Hindi through Persian.
\end{cousinweb}

\begin{grammarlens}[title={What you've built}]
One more word for giving reasons and opinions.
\end{grammarlens}

\subsection*{Your turn}
\begin{itemize}
\raggedright
  \item \emph{Say it:} \emph{qasūr}
  \item \emph{Say it:} \emph{qasūr}, once more
  \item \emph{Say it:} \emph{yah tumhārā qasūr nahīṃ hai.} — It isn't your fault.
  \item \emph{Recall:} say the Hindi for to commend, then the Hindi for praise, then say \emph{qasūr} again
\end{itemize}

\begin{tcolorbox}[breakable,colback=teal!4,colframe=teal!35!black,title={Before you move on}]
What does \hi{क़सूर} mean? (\textquotedblleft{}Fault, blame\textquotedblright{}.) Say it once more.
\end{tcolorbox}
\section[\texorpdfstring{\hi{आलोचना} (ālocnā)}{आलोचना (ālocnā)}]{\hi{आलोचना} (ālocnā) — criticism}

\label{lesson:HI-C367-alocna}



\begin{quote}
Before the new one: say the Hindi for praise, then the Hindi for fault.
\end{quote}

\subsection*{You'll want to know: \hi{आलोचना}}
\textbf{\hi{आलोचना}} — \emph{ālocnā} — \textquotedblleft{}criticism\textquotedblright{}.

\begin{quote}
\hi{सरकार} \hi{की} \hi{नई} \hi{नीति} \hi{की} \hi{बहुत} \hi{आलोचना} \hi{हुई।}
\end{quote}

\begin{quote}
\emph{sarkār kī naī nīti kī bahut ālocnā huī.}
\end{quote}

\begin{quote}
The government's new policy was much criticised.
\end{quote}

Feminine.

\begin{cousinweb}
\textbf{\hi{आलोचना}} (\emph{ālocnā}) is a tatsama, borrowed from Sanskrit \emph{ālocana}, \textquotedblleft{}considering, reflecting\textquotedblright{} (Monier-Williams). Hindi uses it for judging something's faults.
\end{cousinweb}

\begin{grammarlens}[title={What you've built}]
One more word for giving reasons and opinions.
\end{grammarlens}

\subsection*{Your turn}
\begin{itemize}
\raggedright
  \item \emph{Say it:} \emph{ālocnā}
  \item \emph{Say it:} \emph{ālocnā}, once more
  \item \emph{Say it:} \emph{sarkār kī naī nīti kī bahut ālocnā huī.} — The government's new policy was much criticised.
  \item \emph{Recall:} say the Hindi for praise, then the Hindi for fault, then say \emph{ālocnā} again
\end{itemize}

\begin{tcolorbox}[breakable,colback=teal!4,colframe=teal!35!black,title={Before you move on}]
What does \hi{आलोचना} mean? (\textquotedblleft{}Criticism\textquotedblright{}.) Say it once more.
\end{tcolorbox}
\section[\texorpdfstring{\hi{झेलना} (jhelnā)}{झेलना (jhelnā)}]{\hi{झेलना} (jhelnā) — to endure, to put up with}

\label{lesson:HI-C367-jhelna}



\begin{quote}
Before the new one: say the Hindi for fault, then the Hindi for criticism.
\end{quote}

\subsection*{You'll want to know: \hi{झेलना}}
\textbf{\hi{झेलना}} — \emph{jhelnā} — \textquotedblleft{}to endure, to put up with\textquotedblright{}.

\begin{quote}
\hi{उसने} \hi{बहुत} \hi{आलोचना} \hi{झेली।}
\end{quote}

\begin{quote}
\emph{usne bahut ālocnā jhelī.}
\end{quote}

\begin{quote}
He put up with a lot of criticism.
\end{quote}

Past, with \hi{ने} (he or she …): \hi{उसने} \hi{दुख} \hi{झेला}, \emph{usne dukh jhelā}.

\begin{cousinweb}
\textbf{\hi{झेलना}} (\emph{jhelnā}) is inherited from \emph{*jhēlati}, \textquotedblleft{}undergoes\textquotedblright{}, a form CDIAL reconstructs from its descendants (CDIAL 5413).
\end{cousinweb}

\begin{grammarlens}[title={What you've built}]
One more word for giving reasons and opinions.
\end{grammarlens}

\subsection*{Your turn}
\begin{itemize}
\raggedright
  \item \emph{Say it:} \emph{jhelnā}
  \item \emph{Say it:} \emph{jhelnā}, once more
  \item \emph{Say it:} \emph{usne bahut ālocnā jhelī.} — He put up with a lot of criticism.
  \item \emph{Recall:} say the Hindi for fault, then the Hindi for criticism, then say \emph{jhelnā} again
\end{itemize}

\begin{tcolorbox}[breakable,colback=teal!4,colframe=teal!35!black,title={Before you move on}]
What does \hi{झेलना} mean? (\textquotedblleft{}To endure, to put up with\textquotedblright{}.) Say it once more.
\end{tcolorbox}
